% Generated draft: requires mathematical and source review.
\documentclass[11pt]{article}
\usepackage[margin=25mm]{geometry}
\usepackage{fontspec}
\setmainfont{Latin Modern Roman}
\usepackage{amsmath,amssymb,mathtools,unicode-math}
\setmathfont{Latin Modern Math}
\usepackage{enumitem,needspace,xurl,hyperref}
\setlength{\parindent}{0pt}
\setlength{\parskip}{5pt}
\setlength{\emergencystretch}{3em}
\newcommand{\N}{\mathbb{N}}
\newcommand{\Z}{\mathbb{Z}}
\newcommand{\Q}{\mathbb{Q}}
\newcommand{\R}{\mathbb{R}}
\newcommand{\C}{\mathbb{C}}
\newcommand{\F}{\mathbb{F}}
\newcommand{\A}{\mathbb{A}}
\newcommand{\PP}{\mathbb{P}}
\newcommand{\E}{\mathbb{E}}
\newcommand{\eps}{\varepsilon}
\newcommand{\dd}{\,\mathrm{d}}
\begin{document}
\section*{UnsolvedMath ID: 700148}
\textbf{Source problem code:} catalogue entry
\subsection{Definitions and mathematical statement}
For a finite group \(G\) and a positive integer \(n\), put
\[
\mathcal H_n(G)=\{H\leq G:|H|=n\}.
\]
The collection \(\mathcal H_n(G)\) is an \(IE_{\bar n}\)-system precisely when it contains a subgroup that is not normal in \(G\). Thus each qualifying order contributes one system, consisting of all subgroups of that order. Define
\[
i(G)=\left|\left\{n\in\Z_{>0}:\exists H\leq G\text{ with }|H|=n\text{ and }H\not\trianglelefteq G\right\}\right|.
\]
Only divisors of \(|G|\) can contribute, so this number is finite.

A group is \emph{soluble} if its derived series terminates at the trivial subgroup, and \emph{non-soluble} otherwise. A group is \emph{simple} if it is nontrivial and has no nontrivial proper normal subgroup. In the following definitions, every quantified group is finite. For \(k\in\Z_{>0}\), define:
\begin{align*}
S&=\{k\in\Z_{>0}:\forall G,\ i(G)=k\Longrightarrow G\text{ is soluble}\},\\
N&=\{k\in\Z_{>0}:\exists G,\ i(G)=k\text{ and }G\text{ is non-soluble}\},\\
T&=\{k\in\Z_{>0}:\exists G,\ i(G)=k\text{ and }G\text{ is simple}\},\\
C&=\{k\in\Z_{>0}:\nexists G,\ i(G)=k\text{ and }G\text{ is simple}\},\\
A&=\bigl\{k\in T:\nexists G,\ i(G)=k\text{ and }G\text{ is both non-soluble and non-simple}\bigr\}.
\end{align*}
These are, respectively, the sets of soluble, non-soluble, simple, composite, and absolutely simple \emph{group-theoretic numbers}. These adjectives describe the stated group-theoretic properties of \(k\), not its arithmetic primality or compositeness. The source imposes no separate existence requirement in the definitions of \(S\) and \(C\); their universal and nonexistence conditions are retained literally. Including cyclic groups of prime order in the definition of simple causes no issue here: their value of \(i(G)\) is zero, whereas \(k\) is required to be positive.

\textbf{Problem (P. I. Trofimov, 1967; Kourovka Notebook 2.78).} Determine whether \(S\) is finite or infinite, and whether \(A\) is finite or infinite. Does there exist a composite but not soluble group-theoretic number; equivalently, is
\[
C\setminus S\neq\varnothing?
\]
The supplied registry labels this problem open; that label is not an independently verified mathematical or literature assessment.
\subsection{Short English statement}
Count the subgroup orders for which a finite group has at least one non-normal subgroup. Trofimov asks whether finitely or infinitely many positive counts force solubility, whether finitely or infinitely many counts occur in simple groups while occurring in no non-soluble non-simple groups, and whether a count can occur in a non-soluble group but in no simple group.
\subsection{Sources}
\begin{itemize}
\item Local UnsolvedMath registry, permanent ID \(700148\), exact source code \texttt{catalogue entry}, title \emph{Kourovka Notebook Problem 2.78}. Supplied record revision: \texttt{d371ee917ac9647fbb06cd6fc529deec6f3dda82}; record updated \texttt{2026-10-16T14:00:00Z}. The supplied status is \emph{open}.
\item \emph{The Kourovka Notebook (Unsolved Problems in Group Theory)}, original problem 2.78, proposed by P. I. Trofimov in the 2nd issue (1967). The supplied background identifies the 21st issue edition (2026), \texttt{arXiv:1401.0300v48}. The supplied collection description names E. I. Khukhro and V. D. Mazurov as editors. This bibliographic attribution is supplied source data and has not been independently checked.
\end{itemize}
\end{document}
